\subsection{Motivation}
\label{sec:prosd_design_motivation}
As established in Sec.~\ref{sec:introduction} and the benchmark diagnostics
(Sec.~\ref{sec:benchmark_statistics}), the fundamental difficulty of
fixed-viewpoint roadside occupancy is the structural asymmetry between
persistent static layout and transient dynamic participants. Under this
structure, flat one-shot voxel classification is biased toward spatially
dominant static patterns and weakens sparse dynamic cues, while simple loss
reweighting~\cite{FocalLoss,ClassBalancedLoss} only adjusts class weights
without changing how static and dynamic evidence are organized. Accordingly, the
design principle of ProSD-Occ is to organize occupancy reasoning from stable
layout to local residual evidence rather than predicting all voxel categories at
once: the model first establishes a sample-adaptive explanation of the
persistent static layout, then uses it as a reference to recover residual
evidence that the static context cannot sufficiently explain---often
corresponding to sparse yet safety-critical dynamic objects---and finally
recomposes static, dynamic, and free-space evidence into a unified semantic
occupancy field. This paradigm exploits the contextual value of stable layout
while suppressing its dominance over sparse dynamic prediction.
\begin{figure*}[!t]
\centering
\includegraphics[width=\textwidth]{figures/fig08-method-overview.png}
\caption{\textbf{Overall framework of \nameplus{}.}
\nameplus{} formulates infrastructure-side occupancy prediction as a
progressive static-to-dynamic reasoning process rather than one-step voxel
classification. Given camera-only, LiDAR-only, or multi-modal inputs,
Modality-flexible Feature Encoding first produces a unified voxel representation
$\bm{F}$. Static Layout Reasoner then estimates static logits and confidence
$(\bm{Z}_{s}, \bm{C}_{s})$, which provide a soft explanation of persistent
infrastructure layout. Guided by this confidence, Static-guided Residual
Modulation suppresses static-dominant responses and derives dynamic-aware
residual features $\bm{F}_{d}$. Residual Dynamic Predictor further
recovers local dynamic evidence $(\bm{Z}_{d}, \bm{C}_{d})$, and Semantic
Recomposition integrates static, dynamic, and free-space evidence into
the final semantic occupancy logits $\bm{Z}_{\mathrm{final}}$, which are converted by softmax into the final
occupancy prediction.
}
\label{fig:method_overview_refined}
\end{figure*}


\subsection{Overall Architecture}
\label{sec:prosd_overview}

As shown in \Cref{fig:method_overview_refined}, ProSD-Occ organizes
infrastructure-side semantic occupancy prediction as a progressive
static-to-dynamic reasoning framework. The framework decouples modality-specific
feature construction from modality-agnostic occupancy reasoning. Given
camera-only, LiDAR-only, or multi-modal inputs, Modality-flexible Feature
Encoding first converts the observations into a unified voxel representation
$\bm{F}$. In the camera-only and LiDAR-only settings, $\bm{F}$ is instantiated
by the corresponding camera or LiDAR voxel feature. In the multi-modal setting,
camera and LiDAR voxel features are fused before entering the subsequent
reasoning process. This unified representation enables the same
static-to-dynamic reasoning pipeline to be applied across all sensing regimes.

Given $\bm{F}$, ProSD-Occ proceeds from stable layout explanation to dynamic
evidence recovery. Static Layout Reasoner first estimates static logits,
probabilities, and confidence $(\bm{Z}_{s},\bm{P}_{s},\bm{C}_{s})$, where
$\bm{C}_{s}$ provides a sample-adaptive soft explanation of persistent
infrastructure layout. Static-guided Residual Modulation then uses
$\bm{C}_{s}$ to attenuate static-dominant responses while preserving
complementary residual information, producing the dynamic-aware feature
$\bm{F}_{d}$. Based on $\bm{F}_{d}$, Residual Dynamic Predictor estimates
dynamic logits, probabilities, and confidence
$(\bm{Z}_{d},\bm{P}_{d},\bm{C}_{d})$, focusing on local and transient object
evidence. Finally, Semantic Recomposition integrates static, dynamic, and
free-space evidence at the logit level and outputs the final occupancy logits
$\bm{Z}_{\mathrm{final}}$.

The resulting data flow is summarized as
\begin{equation}
  \bm{F}
  \rightarrow
  (\bm{Z}_{s},\bm{P}_{s},\bm{C}_{s})
  \rightarrow
  \bm{F}_{d}
  \rightarrow
  (\bm{Z}_{d},\bm{P}_{d},\bm{C}_{d})
  \rightarrow
  \bm{Z}_{\mathrm{final}},
\end{equation}
where $\bm{Z}_{\mathrm{final}}$ is defined over the unified occupancy label
space and is converted into semantic occupancy probabilities by softmax.

\subsection{Modality-flexible Feature Encoding}
\label{sec:prosd_modality_flexible}
ProSD-Occ performs static-to-dynamic occupancy reasoning on a unified voxel
representation, regardless of the input sensing modality. Given camera images,
LiDAR points, or both, the modality-specific front-end first maps the
observations into a shared 3D voxel feature tensor
\begin{equation}
  \bm{F}\in\mathbb{R}^{B\times C_f\times Z\times H\times W},
\end{equation}
which serves as the common input to the subsequent Progressive
Static-to-Dynamic Reasoning module.

For camera input, we follow common image-to-BEV/voxel encoding
designs~\cite{BEVDet,BEVDepth}. Multi-view images are first processed by an
image backbone and neck, and then lifted into the 3D voxel space through a
depth-aware view transformation. The resulting camera voxel features are
further refined with a BEVFormer-style backward projection module~\cite{BEVFormer}
to aggregate multi-view geometric context. We denote the resulting camera voxel
feature as $\bm{F}_{cam}$.
For LiDAR input, point clouds are converted into voxel representations following
standard voxel-based point cloud encoding~\cite{VoxelNet,PointPillars}. The
voxelized points are encoded by a sparse LiDAR encoder and then processed by a
3D convolutional backbone to obtain LiDAR voxel features, denoted as
$\bm{F}_{lidar}$. This provides a geometry-explicit representation that is
naturally aligned with the occupancy voxel grid.
In the multi-modal setting, we follow the unified BEV/voxel fusion paradigm in
multi-sensor perception~\cite{BEVFusion,OpenOccupancy}. Camera and LiDAR
features are first projected into the same voxel grid and then fused before
occupancy reasoning. Concretely, the fusion block computes lightweight
attention maps from each modality and uses them to cross-modulate the voxel
features of the other modality. The modulated camera and LiDAR features are
concatenated and projected back to the shared channel dimension by a 3D
convolutional block, producing the fused voxel feature
$\Phi(\bm{F}_{cam},\bm{F}_{lidar})$.
Therefore, the input to ProSD-Occ is uniformly defined as
\begin{equation}
  \bm{F} =
  \begin{cases}
  \bm{F}_{cam}, & \text{Camera-only},\\
  \bm{F}_{lidar}, & \text{LiDAR-only},\\
  \Phi(\bm{F}_{cam},\bm{F}_{lidar}), & \text{Multi-modal}.
  \end{cases}
\end{equation}
This formulation decouples modality-dependent feature construction from the
subsequent modality-agnostic static-to-dynamic occupancy reasoning, allowing
the same ProSD-Occ framework to be applied to camera-only, LiDAR-only, and
multi-modal infrastructure perception.

\subsection{Progressive Static-to-Dynamic Reasoning}
\label{sec:prosd_reasoning}
Following the overview in Sec.~\ref{sec:prosd_overview}, we now detail the
four components of the progressive reasoning module
(\Cref{fig:method_overview_refined}): the Static Layout Reasoner,
Static-guided Residual Modulation, the Residual Dynamic Predictor, and Semantic
Recomposition.
\begin{table*}[!t]
\caption{\textbf{Main InfraOcc occupancy results.} Every row, including
RoadOcc-S, uses strictly current-frame-only input; no history-enabled method
appears here. Temporal variants are evaluated separately in the joint
occupancy--velocity comparison. Rankings are computed within each
sensing regime; best, second-best, and third-best results are highlighted.}
\label{tab:infraocc_main_results}
\centering
\definecolor{best}{RGB}{255, 180, 180}
\definecolor{second}{RGB}{255, 230, 180}
\definecolor{third}{RGB}{255, 255, 200}
\newlength{\mainmethodwd}
\newlength{\mainmetricwd}
\newlength{\mainclasswd}
\setlength{\mainmethodwd}{0.135\textwidth}
\setlength{\mainmetricwd}{0.061\textwidth}
\setlength{\mainclasswd}{\dimexpr(\textwidth-\mainmethodwd-4\mainmetricwd-2\arrayrulewidth)/15\relax}
\newcommand{\methodcell}[1]{\makebox[\mainmethodwd][l]{#1}}
\newcommand{\metriccell}[1]{\makebox[\mainmetricwd][c]{#1}}
\newcommand{\classcell}[1]{\makebox[\mainclasswd][c]{#1}}
\newcommand{\bestmetriccell}[1]{\cellcolor{best}\metriccell{\textbf{#1}}}
\newcommand{\secondmetriccell}[1]{\cellcolor{second}\metriccell{\underline{#1}}}
\newcommand{\thirdmetriccell}[1]{\cellcolor{third}\metriccell{\emph{#1}}}
\newcommand{\bestclasscell}[1]{\cellcolor{best}\classcell{\textbf{#1}}}
\newcommand{\secondclasscell}[1]{\cellcolor{second}\classcell{\underline{#1}}}
\newcommand{\thirdclasscell}[1]{\cellcolor{third}\classcell{\emph{#1}}}
\scriptsize
\setlength{\tabcolsep}{0pt}
\renewcommand{\arraystretch}{1.05}
\makebox[\textwidth][c]{%
\begin{tabular}{@{}c|cccc|*{15}{c}@{}}
\toprule
\methodcell{Method} & \metriccell{gIoU} & \metriccell{$\mathrm{mIoU}_{\mathrm{all}}$} & \metriccell{$\mathrm{mIoU}_{\mathrm{dyn}}$} & \metriccell{$\mathrm{mIoU}_{\mathrm{sta}}$} & \classcell{\occclassname{Bike}} & \classcell{\occclassname{Bus}} & \classcell{\occclassname{Car}} & \classcell{\occclassname{Motorcycle}} & \classcell{\occclassname{Pedestrian}} & \classcell{\occclassname{Truck}} & \classcell{\occclassname{Other}} & \classcell{\occclassname{Barrier}} & \classcell{\occclassname{Cone}} & \classcell{\occclassname{Driveable}} & \classcell{\occclassname{Sidewalk}} & \classcell{\occclassname{Terrain}} & \classcell{\occclassname{Manmade}} & \classcell{\occclassname{Vegetation}} & \classcell{\occclassname{Free}} \\
\methodcell{} & \metriccell{} & \metriccell{} & \metriccell{} & \metriccell{} & \classcell{\classswatch{occBike}} & \classcell{\classswatch{occBus}} & \classcell{\classswatch{occCar}} & \classcell{\classswatch{occMotor}} & \classcell{\classswatch{occPed}} & \classcell{\classswatch{occTruck}} & \classcell{\classswatch{occOther}} & \classcell{\classswatch{occBarrier}} & \classcell{\classswatch{occCone}} & \classcell{\classswatch{occDrive}} & \classcell{\classswatch{occSide}} & \classcell{\classswatch{occTerrain}} & \classcell{\classswatch{occManmade}} & \classcell{\classswatch{occVegetation}} & \classcell{\classswatch{occFree}} \\
\midrule
\rowcolor{gray!5}{\methodcell{\textit{Multi-modal (C+L)}}}  \\
% M-CONet~\cite{OpenOccupancy} & 81.40 & 57.23 & 34.78 & 74.06 & 22.52 & 65.49 & 54.20 & 19.95 & 32.26 & 14.29 & 48.78 & 76.79 & 66.67 & 84.64 & 80.74 & 84.18 & 76.73 & 73.99 & 97.51 \\
\methodcell{BEVFusion~\cite{BEVFusion}} & \metriccell{\secondcell{77.01}} & \metriccell{\secondcell{55.44}} & \metriccell{\secondcell{32.57}} & \metriccell{\secondcell{72.60}} & \classcell{\secondcell{11.03}} & \classcell{\secondcell{65.04}} & \classcell{\secondcell{55.73}} & \classcell{\secondcell{15.26}} & \classcell{\secondcell{29.48}} & \classcell{\bestcell{18.88}} & \classcell{\secondcell{55.04}} & \classcell{\secondcell{72.04}} & \classcell{\secondcell{72.07}} & \classcell{\secondcell{75.43}} & \classcell{\secondcell{73.53}} & \classcell{\secondcell{81.26}} & \classcell{\secondcell{76.82}} & \classcell{\secondcell{74.60}} & \classcell{\secondcell{96.58}} \\
% BEVDepth~\cite{BEVDepth} & 76.20 & 55.34 & 33.07 & 72.04 & 12.61 & 65.35 & 55.37 & 17.51 & 29.61 & 17.95 & 54.96 & 73.50 & 71.20 & 74.79 & 73.09 & 79.69 & 75.53 & 73.56 & 96.41 \\
\rowcolor{gray!5}\methodcell{ProSD-Occ} & \metriccell{\bestcell{91.78}} & \metriccell{\bestcell{65.87}} & \metriccell{\bestcell{37.92}} & \metriccell{\bestcell{86.82}} & \classcell{\bestcell{30.07}} & \classcell{\bestcell{66.95}} & \classcell{\bestcell{58.55}} & \classcell{\bestcell{19.82}} & \classcell{\bestcell{34.65}} & \classcell{\secondcell{17.51}} & \classcell{\bestcell{72.31}} & \classcell{\bestcell{73.79}} & \classcell{\bestcell{89.68}} & \classcell{\bestcell{94.25}} & \classcell{\bestcell{88.35}} & \classcell{\bestcell{94.73}} & \classcell{\bestcell{90.37}} & \classcell{\bestcell{91.12}} & \classcell{\bestcell{99.09}} \\
\midrule
\rowcolor{gray!5}{\methodcell{\textit{LiDAR-Only (L)}}}  \\
% L-CONet~\cite{OpenOccupancy} & 83.25 & 57.76 & 32.48 & 76.72 & 10.67 & 70.74 & 57.16 & 14.87 & 28.81 & 12.65 & 53.46 & 76.57 & 72.23 & 85.55 & 82.77 & 86.40 & 79.90 & 76.85 & 97.76 \\
\methodcell{VoxelNet~\cite{VoxelNet}} & \metriccell{\secondcell{75.69}} & \metriccell{\secondcell{53.87}} & \metriccell{31.13} & \metriccell{\secondcell{70.92}} & \classcell{11.08} & \classcell{\bestcell{65.63}} & \classcell{\secondcell{54.52}} & \classcell{\secondcell{13.43}} & \classcell{27.85} & \classcell{14.29} & \classcell{\secondcell{53.39}} & \classcell{72.22} & \classcell{\secondcell{68.26}} & \classcell{\secondcell{75.13}} & \classcell{\secondcell{72.35}} & \classcell{\secondcell{78.99}} & \classcell{\secondcell{74.69}} & \classcell{\secondcell{72.36}} & \classcell{\secondcell{96.31}} \\
\methodcell{PointPillars~\cite{PointPillars}} & \metriccell{70.53} & \metriccell{50.51} & \metriccell{\secondcell{31.55}} & \metriccell{64.74} & \classcell{\bestcell{14.99}} & \classcell{61.62} & \classcell{54.03} & \classcell{11.50} & \classcell{\secondcell{32.62}} & \classcell{\secondcell{14.55}} & \classcell{44.06} & \classcell{\bestcell{73.83}} & \classcell{57.49} & \classcell{72.11} & \classcell{67.67} & \classcell{70.75} & \classcell{66.42} & \classcell{65.56} & \classcell{95.27} \\
\rowcolor{gray!5}\methodcell{ProSD-Occ} & \metriccell{\bestcell{90.86}} & \metriccell{\bestcell{63.66}} & \metriccell{\bestcell{33.87}} & \metriccell{\bestcell{85.99}} & \classcell{\secondcell{11.59}} & \classcell{\secondcell{65.62}} & \classcell{\bestcell{58.29}} & \classcell{\bestcell{15.22}} & \classcell{\bestcell{33.09}} & \classcell{\bestcell{19.42}} & \classcell{\bestcell{68.00}} & \classcell{\secondcell{73.75}} & \classcell{\bestcell{90.18}} & \classcell{\bestcell{94.02}} & \classcell{\bestcell{88.14}} & \classcell{\bestcell{95.16}} & \classcell{\bestcell{88.53}} & \classcell{\bestcell{90.17}} & \classcell{\bestcell{98.99}} \\
\midrule
\rowcolor{gray!5}{\methodcell{\textit{Camera-Only (C)}}} \\
\methodcell{BEVDet~\cite{BEVDet}} & \metriccell{68.62} & \metriccell{45.14} & \thirdmetriccell{21.19} & \metriccell{63.11} & \classcell{9.23} & \thirdclasscell{42.75} & \classcell{38.04} & \thirdclasscell{9.68} & \classcell{15.01} & \secondclasscell{12.45} & \classcell{42.88} & \classcell{70.36} & \classcell{53.77} & \classcell{72.04} & \classcell{68.07} & \classcell{73.09} & \classcell{64.28} & \classcell{60.38} & \classcell{94.87} \\
\methodcell{BEVFormer~\cite{BEVFormer}} & \thirdmetriccell{81.72} & \metriccell{47.76} & \metriccell{10.73} & \thirdmetriccell{75.53} & \classcell{7.97} & \classcell{20.39} & \classcell{17.65} & \classcell{3.47} & \classcell{8.62} & \classcell{6.30} & \classcell{42.47} & \thirdclasscell{76.26} & \secondclasscell{72.83} & \classcell{85.19} & \thirdclasscell{81.95} & \thirdclasscell{88.95} & \thirdclasscell{79.52} & \thirdclasscell{77.05} & \thirdclasscell{97.56} \\
\methodcell{BEVDepth~\cite{BEVDepth}} & \metriccell{71.89} & \metriccell{46.81} & \metriccell{19.64} & \metriccell{67.19} & \secondclasscell{14.84} & \classcell{33.87} & \classcell{38.74} & \classcell{3.42} & \classcell{16.86} & \classcell{10.14} & \classcell{48.73} & \classcell{71.05} & \classcell{62.37} & \classcell{73.31} & \classcell{69.32} & \classcell{76.07} & \classcell{70.51} & \classcell{66.15} & \classcell{95.54} \\
\methodcell{TPVFormer~\cite{TPVFormer}} & \metriccell{73.80} & \metriccell{43.09} & \metriccell{4.36} & \metriccell{72.13} & \classcell{0.10} & \classcell{3.69} & \classcell{13.80} & \classcell{0.41} & \classcell{6.86} & \classcell{1.33} & \classcell{42.86} & \classcell{75.29} & \classcell{69.69} & \classcell{80.61} & \classcell{77.62} & \classcell{86.63} & \classcell{72.61} & \classcell{71.70} & \classcell{96.12} \\
\methodcell{SurroundOcc~\cite{SurroundOcc}} & \secondmetriccell{83.78} & \secondmetriccell{50.77} & \metriccell{11.37} & \secondmetriccell{80.32} & \classcell{12.66} & \classcell{21.90} & \classcell{18.54} & \classcell{2.06} & \classcell{9.50} & \classcell{3.55} & \thirdclasscell{65.09} & \bestclasscell{79.04} & \thirdclasscell{72.70} & \thirdclasscell{86.38} & \secondclasscell{83.99} & \secondclasscell{91.10} & \secondclasscell{83.16} & \secondclasscell{81.09} & \secondclasscell{97.80} \\
\methodcell{CONet~\cite{OpenOccupancy}} & \metriccell{78.56} & \metriccell{49.25} & \metriccell{18.93} & \metriccell{72.00} & \classcell{3.06} & \classcell{37.24} & \thirdclasscell{39.47} & \classcell{4.42} & \thirdclasscell{17.64} & \classcell{11.72} & \classcell{45.01} & \classcell{75.31} & \classcell{64.94} & \classcell{84.29} & \classcell{79.45} & \classcell{84.27} & \classcell{72.96} & \classcell{69.79} & \classcell{97.08} \\
\methodcell{STCOcc~\cite{STCOcc}} & \metriccell{78.54} & \thirdmetriccell{50.71} & \secondmetriccell{23.48} & \metriccell{71.14} & \thirdclasscell{13.43} & \secondclasscell{47.83} & \secondclasscell{39.62} & \secondclasscell{10.21} & \secondclasscell{17.83} & \thirdclasscell{11.93} & \classcell{50.18} & \classcell{71.54} & \classcell{58.33} & \secondclasscell{90.30} & \classcell{78.78} & \classcell{84.90} & \classcell{67.72} & \classcell{67.32} & \classcell{97.47} \\
\methodcell{SparseOcc~\cite{SparseOcc}} & \metriccell{71.61} & \metriccell{45.42} & \metriccell{9.41} & \metriccell{72.42} & \classcell{1.50} & \classcell{20.94} & \classcell{24.87} & \classcell{0.00} & \classcell{4.60} & \classcell{4.57} & \bestclasscell{67.68} & \secondclasscell{76.37} & \classcell{61.51} & \classcell{74.67} & \classcell{76.59} & \classcell{88.54} & \classcell{68.79} & \classcell{65.19} & \classcell{96.31} \\
\methodcell{GaussianFormer~\cite{GaussianFormer}} & \metriccell{77.89} & \metriccell{43.63} & \metriccell{10.45} & \metriccell{68.51} & \classcell{3.13} & \classcell{22.11} & \classcell{23.17} & \classcell{1.29} & \classcell{8.21} & \classcell{4.80} & \classcell{36.70} & \classcell{68.54} & \classcell{55.60} & \classcell{82.59} & \classcell{79.37} & \classcell{74.20} & \classcell{77.85} & \classcell{73.21} & \classcell{96.71} \\
\rowcolor{gray!5}\methodcell{ProSD-Occ} & \bestmetriccell{89.08} & \bestmetriccell{60.08} & \bestmetriccell{28.97} & \bestmetriccell{83.42} & \bestclasscell{22.74} & \bestclasscell{55.45} & \bestclasscell{47.84} & \bestclasscell{12.27} & \bestclasscell{21.01} & \bestclasscell{14.49} & \secondclasscell{65.99} & \classcell{72.72} & \bestclasscell{82.14} & \bestclasscell{93.84} & \bestclasscell{86.88} & \bestclasscell{93.31} & \bestclasscell{86.09} & \bestclasscell{86.38} & \bestclasscell{98.78} \\
\bottomrule
\end{tabular}%
}
\end{table*}

\subsubsection{Static Layout Reasoner.}
\label{sec:prosd_static_reasoner}
This module estimates persistent static layout from the unified voxel feature.
Its goal is to explain stable infrastructure elements before dynamic reasoning
is performed. Let $\mathcal{S}$ denote the static class set, including road
surfaces, sidewalks, buildings, vegetation, barriers, and other persistent
infrastructure categories. As shown in \Cref{fig:method_overview_refined}, the
Static Layout Reasoner is implemented by a compact BasicBlock3D tower followed
by a voxel-wise MLP predictor. Each BasicBlock3D is a residual 3D convolutional
block composed of two $3\times3\times3$ Conv3D-BN layers and a skip connection,
with ReLU activation applied after residual aggregation. The MLP predictor is
shared across voxels and consists of two linear layers with a Softplus
activation, mapping each voxel feature vector to logits over $\mathcal{S}$ and
an additional ``other'' class:
\begin{equation}
  \bm{Z}_{s}=h_s(\bm{F}),
  \quad
  \bm{Z}_{s}\in
  \mathbb{R}^{B\times(|\mathcal{S}|+1)\times Z\times H\times W} .
\end{equation}

Different from a generic background category, the ``other'' class in the static
branch explicitly represents voxels that are not explained by static layout,
including dynamic objects and free space. Therefore, the static branch is not
asked to solve the full semantic occupancy problem. Instead, it only estimates
whether each voxel can be explained by persistent static structure. The static
probability is computed by
\begin{equation}
  \bm{P}_{s}=\operatorname{softmax}(\bm{Z}_{s}).
\end{equation}
We define the static confidence as the probability that a voxel does not belong
to the static-branch ``other'' class:
\begin{equation}
  \bm{C}_{s}
  =
  1-\bm{P}_{s}(\mathrm{other}).
\end{equation}
The confidence $\bm{C}_{s}$ provides a sample-adaptive soft explanation of
persistent layout. It is not used as a hard static mask; instead, it guides the
subsequent residual modulation step, where static-dominant responses are
suppressed to expose potential dynamic evidence.


\subsubsection{Static-guided Residual Modulation.}
\label{sec:prosd_residual_modulation}

Static-guided Residual Modulation bridges static layout reasoning and residual
dynamic prediction by transforming the unified voxel feature $\bm{F}$ into a
dynamic-aware feature $\bm{F}_{d}$. As shown in
\Cref{fig:method_overview_refined}, this module contains two complementary
paths: a suppression path that attenuates static-dominant responses according
to the predicted static confidence, and a residual path that preserves raw voxel
evidence to avoid over-suppression.

Given the static confidence $\bm{C}_{s}$, the suppression path constructs a
voxel-wise suppression gate:
\begin{equation}
  \alpha=\sigma(a),
  \quad
  \bm{G}
  =
  \operatorname{clamp}
  \left(1-\alpha\,\operatorname{sg}(\bm{C}_{s}),0,1\right),
\end{equation}
where $a$ is a learnable scalar, $\alpha\in(0,1)$ controls the suppression
strength, and $\operatorname{sg}(\cdot)$ denotes stop-gradient. Voxels with high
static confidence are softly attenuated, encouraging the subsequent dynamic
branch to focus on evidence not sufficiently explained by static layout. The
static-suppressed feature is obtained by
\begin{equation}
  \bm{F}_{sup}=\bm{F}\odot\bm{G}.
\end{equation}

Since static confidence can be uncertain around object boundaries, occluded
regions, or sparse dynamic objects, pure suppression may discard useful dynamic
cues. The residual path therefore introduces a bypass from the original voxel
feature:
\begin{equation}
  \bm{F}_{raw}
  =
  \operatorname{BN}
  \left(\operatorname{Conv}_{1\times1\times1}(\bm{F})\right),
  \quad
  \beta=\sigma(b),
\end{equation}
where $b$ is a learnable scalar and $\beta\in(0,1)$ controls the bypass
strength. The final dynamic-aware feature is then formed as
\begin{equation}
  \bm{F}_{d}
  =
  \bm{F}_{sup}+\beta\bm{F}_{raw}.
\end{equation}
In this way, the module suppresses dominant static responses while retaining
complementary residual evidence, providing a more suitable representation for
dynamic occupancy prediction.

\subsubsection{Residual Dynamic Predictor.}
\label{sec:prosd_dynamic_predictor}
The Residual Dynamic Predictor estimates transient traffic participants from
the dynamic-aware feature $\bm{F}_{d}$. Different from the Static Layout
Reasoner, this module focuses on local, sparse, and boundary-sensitive object
evidence that remains after static-guided modulation. It adopts the same
compact BasicBlock3D tower and voxel-wise MLP predictor, but predicts over the
dynamic class set $\mathcal{D}$ and an additional ``other'' class:
\begin{equation}
  \bm{Z}_{d}=h_d(\bm{F}_{d}),
  \quad
  \bm{Z}_{d}\in
  \mathbb{R}^{B\times(|\mathcal{D}|+1)\times Z\times H\times W} .
\end{equation}
Here, the dynamic-branch ``other'' class denotes voxels that are not explained
by dynamic objects, including static structures and free space. The dynamic
probability and confidence are computed as
\begin{equation}
  \bm{P}_{d}=\operatorname{softmax}(\bm{Z}_{d}),
  \quad
  \bm{C}_{d}=1-\bm{P}_{d}(\mathrm{other}).
\end{equation}
Since $\bm{F}_{d}$ has already been modulated by static confidence, this
branch is encouraged to recover residual dynamic evidence rather than relearn
dominant static layout from scratch.

% Queue the main double-column table before Sec. 5 so it can appear earlier.

\begin{figure*}[!t]
  \centering
  \includegraphics[width=0.98\textwidth]{figures/fig09-modality-qualitative.png}
  \caption{\textbf{Qualitative comparison.}
  The two rows show validation samples with aligned roadside camera evidence,
  ground-truth occupancy, and camera-only occupancy predictions from TPVFormer,
  SparseOCC, STCOcc, and ProSD-Occ (C). STCOcc denotes the plain camera
  baseline, while ProSD-Occ (C) applies progressive static-to-dynamic
  reasoning.}
  \label{fig:infraocc_modality_qualitative}
\end{figure*}

\subsubsection{Semantic Recomposition.}
\label{sec:prosd_semantic_recomposition}
After static and dynamic reasoning, the two branches provide complementary
partial explanations of the occupancy field. The static branch is specialized
for persistent infrastructure layout, while the dynamic branch focuses on
transient traffic participants. However, neither branch alone is defined over
the complete semantic label space. A direct averaging or concatenation of their
outputs would ignore their branch specialization and blur the distinction
between static evidence, dynamic evidence, and free-space evidence. Therefore,
we recompose the branch outputs at the logit level.
Specifically, given the static outputs $(\bm{P}_{s},\bm{C}_{s})$ and dynamic
outputs $(\bm{P}_{d},\bm{C}_{d})$, we construct a voxel-wise recomposition
input:
\begin{equation}
  \bm{x}_{g}
  =
  [\bm{P}_{s},\bm{C}_{s},\bm{P}_{d},\bm{C}_{d}].
\end{equation}
A lightweight MLP followed by a sigmoid function predicts three group-wise
recomposition weights:
\begin{equation}
  [\bm{w}_{\mathrm{static}},\bm{w}_{\mathrm{dynamic}},\bm{w}_{\mathrm{free}}]
  =
  \sigma(g_{\phi}(\bm{x}_{g})),
\end{equation}
where $\bm{w}_{\mathrm{static}}$, $\bm{w}_{\mathrm{dynamic}}$, and
$\bm{w}_{\mathrm{free}}$ correspond to
the static classes, dynamic classes, and free-space class, respectively.

The branch logits $\bm{Z}_{s}$ and $\bm{Z}_{d}$ are then recomposed according to
their semantic groups. Let $\bm{Z}_{s,o}$ and $\bm{Z}_{d,o}$ denote the
``other'' logits of the static and dynamic branches. For a static class
$c\in\mathcal{S}$, the final logit is composed from the corresponding static
logit and the dynamic-branch ``other'' logit:
\begin{equation}
  \bm{Z}_{\mathrm{final},c}
  =
  \bm{w}_{\mathrm{static}}\bm{Z}_{s,c}
  +(1-\bm{w}_{\mathrm{static}})\bm{Z}_{d,o}.
\end{equation}
For a dynamic class $c\in\mathcal{D}$, the final logit is composed from the
dynamic logit and the static-branch ``other'' logit:
\begin{equation}
  \bm{Z}_{\mathrm{final},c}
  =
  \bm{w}_{\mathrm{dynamic}}\bm{Z}_{s,o}
  +(1-\bm{w}_{\mathrm{dynamic}})\bm{Z}_{d,c}.
\end{equation}
For the free-space class $f$, the final logit is obtained by combining the
``other'' logits from both branches:
\begin{equation}
  \bm{Z}_{\mathrm{final},f}
  =
  \bm{w}_{\mathrm{free}}\bm{Z}_{s,o}
  +(1-\bm{w}_{\mathrm{free}})\bm{Z}_{d,o}.
\end{equation}

The final semantic occupancy probability is computed as
\begin{equation}
  \bm{P}_{\mathrm{final}}
  =
  \operatorname{softmax}(\bm{Z}_{\mathrm{final}}).
\end{equation}
During training, $\bm{Z}_{\mathrm{final}}$ is directly passed to the occupancy
loss, while $\bm{P}_{\mathrm{final}}$ is used for semantic occupancy prediction.

% Queue the main qualitative figure before the ablation floats.


\subsection{Optimization Objectives}
\label{sec:prosd_objective}

We supervise three outputs in ProSD-Occ: the final occupancy prediction and the
two intermediate static and dynamic branch predictions. The final logits
$\bm{Z}_{\mathrm{final}}$ are supervised in the complete semantic occupancy
space, yielding $\mathcal{L}_{\mathrm{occ}}$. For the static and dynamic
branches, semantic classes outside the corresponding group are mapped to the
branch-specific ``other'' class, yielding $\mathcal{L}_{s}$ and
$\mathcal{L}_{d}$. All three losses adopt the same composite occupancy
objective:
\begin{equation}
  \mathcal{L}_{r}
  =
  \mathcal{L}_{r}^{\mathrm{CE}}
  +
  \mathcal{L}_{r}^{\mathrm{Sem}}
  +
  \mathcal{L}_{r}^{\mathrm{Lovasz}},
  \quad r\in\{\mathrm{occ},s,d\},
\end{equation}
where $\mathcal{L}^{\mathrm{CE}}$, $\mathcal{L}^{\mathrm{Sem}}$, and
$\mathcal{L}^{\mathrm{Lovasz}}$ denote cross-entropy loss, semantic scaling
loss~\cite{MonoScene}, and Lovasz loss~\cite{Berman_2018_CVPR}, respectively.

Finally, since $\bm{C}_{s}$ directly controls the static-guided suppression
gate, its stability matters: persistent structures observed from a fixed
viewpoint should yield consistent static confidence across same-scene frames,
yet per-sample supervision alone leaves it fluctuating under transient
occlusions, dynamic objects, and sensing noise. We therefore add a
static-confidence consistency regularizer on high-confidence static voxels.
Because all frames are annotated in the same fixed roadside coordinate
system, the static confidences of frames from the same scene are voxel-aligned
and directly comparable. Within each minibatch, let $\mathcal{B}$ be the set of
samples from the same scene; we define the per-voxel consensus as the
stop-gradient batch mean
\begin{equation}
  \bar{\bm{C}}_{s}(v)
  =
  \operatorname{sg}\!\Big(
  \tfrac{1}{|\mathcal{B}|}\textstyle\sum_{i\in\mathcal{B}}\bm{C}_{s,i}(v)
  \Big),
\end{equation}
and restrict the regularizer to the high-confidence static set
$\Omega=\{v:\bar{\bm{C}}_{s}(v)>\tau\}$ with threshold $\tau$. The consistency
loss then penalizes each frame's deviation from this consensus,
\begin{equation}
  \mathcal{L}_{cons}
  =
  \frac{1}{|\Omega|}
  \sum_{i,\,v\in\Omega}
  \left(
  \bm{C}_{s,i}(v)-\bar{\bm{C}}_{s}(v)
  \right)^2 .
\end{equation}
The stop-gradient makes the consensus a stable target and prevents the
trivial collapse of all confidences to a common value, so the regularizer only
removes per-frame fluctuations on confidently static voxels.

The overall objective is
\begin{equation}
  \mathcal{L}
  =
  \mathcal{L}_{occ}
  +
  \lambda_s\mathcal{L}_{s}
  +
  \lambda_d\mathcal{L}_{d}
  +
  \lambda_c\mathcal{L}_{cons}.
\end{equation}

